\chapter{1977年新疆卷（文）}

\section{解答题}

\begin{problem}
解方程：\( \frac{x-1}{x+1}-\frac{5}{1-x}=\frac{4}{x^2-1} \).
\end{problem}

\begin{solution}
由分母不为零，得 \(x\ne -1\)，\(x\ne 1\)，所以原方程的定义域为 \(x\ne\pm1\). 在此条件下，两边同乘 \(x^2-1\)，得 \(\left(x-1\right)^2+5\left(x+1\right)=4\)，即 \(x^2-2x+1+5x+5=4\)，整理得 \(x^2+3x+2=0\)，所以 \(\left(x+1\right)\left(x+2\right)=0\)，从而 \(x=-1\) 或 \(x=-2\).

其中 \(x=-1\) 不在原方程的定义域内，应舍去. 代入检验 \(x=-2\)：左边为 \(3-\frac{5}{3}=\frac{4}{3}\)，右边为 \(\frac{4}{3}\)，等式成立. 因此原方程的解为 \(x=-2\).
\end{solution}

\begin{problem}
红心大队贫下中农坚决顶住“四人帮”对农业学大寨运动的干扰和破坏，积极发展养猪事业，1974年底全大队养猪\(200\text{ 头}\)，到1976年底发展到\(1200\text{ 头}\)，求平均每年的增长率，（精确到\(1\%\)）.
\end{problem}

\begin{solution}
设平均每年的增长率为 \(x\). 增长率不能为负，因此 \(x\geq0\). 从 1974 年底到 1976 年底经过两年，按每年增长率均为 \(x\) 计算，猪数满足 \(1200=200\left(1+x\right)^2\)，约去 \(200\) 得 \(\left(1+x\right)^2=6\). 由于 \(1+x>0\)，所以只能取正平方根，得到 \(1+x=\sqrt{6}\)，即 \(x=\sqrt{6}-1\).

计算得 \(\sqrt{6}\approx2.44949\)，故 \(x\approx1.44949=144.949\%\). 精确到 \(1\%\) 为 \(145\%\).

答：平均每年的增长率为 \(145\%\).
\end{solution}

\begin{problem}
计算：

\begin{enumerate}
\item 化简：\( \frac{\frac{3}{4}\left( ab \right)^{-\frac{1}{3}}\left( \frac{4}{9}a^4b^{\frac{2}{3}} \right)^{\frac{1}{2}}}{2a^3\sqrt{a^2}} \).

\item 已知 \( \lg 2=0.3010 \)，求 \( \lg \sqrt{800} \).
\end{enumerate}
\end{problem}

\begin{solution}
\begin{enumerate}
\item 按题目中分数指数幂的通常定义，取 \(a>0\)，\(b>0\)，于是 \(\sqrt{a^2}=a\). 利用幂的运算性质，\(\left(ab\right)^{-\frac{1}{3}}=a^{-\frac{1}{3}}b^{-\frac{1}{3}}\)，且 \(\left(\frac{4}{9}a^4b^{\frac{2}{3}}\right)^{\frac{1}{2}}=\frac{2}{3}a^2b^{\frac{1}{3}}\). 因此原式的系数为 \(\frac{3}{4}\times\frac{2}{3}\div2=\frac14\)，\(a\) 的指数为 \(-\frac13+2-3-1=-\frac73\)，\(b\) 的指数为 \(-\frac13+\frac13=0\)，从而原式为 \(\frac{1}{4}a^{-\frac73}=\frac{1}{4a^{\frac73}}=\frac{\sqrt[3]{a^2}}{4a^3}\).

\item 因为 \(800=2^3\times10^2\)，所以 \(\lg\sqrt{800}=\frac{1}{2}\lg800=\frac{1}{2}\lg\left(2^3\times10^2\right)=\frac{3}{2}\lg2+\lg10=\frac{3}{2}\times0.3010+1=1.4515\).
\end{enumerate}
\end{solution}

\begin{problem}
求证：四边形对边中点的连线互相平分.

\bitmapfigure[max width=0.36\linewidth]{img/1977/xinjiang_liberal/q4_fig1.png}
\end{problem}

\begin{solution}
连接 \(BD\). 因为 \(M\)，\(N\) 分别是 \(BC\)，\(CD\) 的中点，所以在 \(\triangle BCD\) 中，依据中位线定理有 \(MN\parallel BD\)，且 \(MN=\frac{1}{2}BD\). 因为 \(F\)，\(E\) 分别是 \(AB\)，\(AD\) 的中点，所以在 \(\triangle BAD\) 中，依据中位线定理有 \(FE\parallel BD\)，且 \(FE=\frac{1}{2}BD\). 于是 \(MN\parallel FE\)，且 \(MN=FE\). 四边形 \(FMNE\) 的一组对边 \(MN\) 和 \(FE\) 平行且相等，因此 \(FMNE\) 是平行四边形. 平行四边形的对角线 \(FN\) 与 \(ME\) 互相平分，所以 \(EM\) 与 \(FN\) 互相平分，命题得证.
\end{solution}

\begin{problem}
解方程组 \( \begin{cases}
2x-3y=9\\
3x+4y=5
\end{cases} \)
\end{problem}

\begin{solution}
将第一个方程乘以 \(3\)，得到 \(6x-9y=27\)；将第二个方程乘以 \(2\)，得到 \(6x+8y=10\). 两式相减，得 \(-17y=17\)，所以 \(y=-1\). 代入第一个方程，得 \(2x-3\left(-1\right)=9\)，即 \(2x=6\)，所以 \(x=3\).

回代检验：\(2\times3-3\times\left(-1\right)=9\)，\(3\times3+4\times\left(-1\right)=5\)，两式均成立. 因此方程组的解为 \(x=3\)，\(y=-1\).
\end{solution}

\begin{problem}
红升人民公社小水库拦水坝的横断面为梯形，坝顶宽\(2.5\text{ 米}\)，坝高\(4\text{ 米}\)，迎水坡，背水坡与水平面的夹角分别是 \(30^\circ\) 和 \(45^\circ\)，求坝底的宽（精确到\(0.1\text{ 米}\)）.
\end{problem}

\begin{solution}
分别由 \(A\)，\(D\) 作 \(AE\perp BC\)，\(DF\perp BC\). 则 \(AE\) 和 \(DF\) 都是梯形的高，所以 \(AE=DF=4\text{ 米}\).

\solutionfigure{\bitmapfigure[max width=0.64\linewidth]{img/1977/xinjiang_liberal/q6_fig1.png}}

在直角三角形 \(ABE\) 中，\(\angle ABE=30^\circ\)，故 \(\tan30^\circ=\frac{AE}{BE}\)，所以 \(BE=\frac{4}{\tan30^\circ}=4\sqrt{3}\text{ 米}\). 在直角三角形 \(DCF\) 中，\(\angle DCF=45^\circ\)，故 \(\tan45^\circ=\frac{DF}{CF}=1\)，所以 \(CF=DF=4\text{ 米}\). 又因为 \(AE\perp BC\)，\(DF\perp BC\)，所以 \(AE\parallel DF\)；结合 \(AD\parallel BC\) 及 \(EF\subset BC\)，可知四边形 \(AEFD\) 是矩形，从而 \(EF=AD=2.5\text{ 米}\). 于是 \(BC=BE+EF+FC=4\sqrt{3}+2.5+4\approx13.428\text{ 米}\).

精确到 \(0.1\text{ 米}\)，得 \(BC\approx13.4\text{ 米}\).

答：坝底宽为 \(13.4\text{ 米}\).
\end{solution}

\begin{problem}
已知直线过 \( A\left( 2,-5 \right) \) 和 \( B\left( -3,0 \right) \) 两点.

\begin{enumerate}
\item 求此直线方程；
\item 指出直线的斜率 \(K\)；
\item 求直线与 \(OX\) 轴正方向的夹角 \(\alpha\).
\end{enumerate}

\bitmapfigure[float=inline,max width=0.38\linewidth]{img/1977/xinjiang_liberal/q7_fig1.png}
\end{problem}

\begin{solution}
\begin{enumerate}
\item 由两点式，直线满足 \(\frac{y-0}{-5-0}=\frac{x-\left(-3\right)}{2-\left(-3\right)}\)，即 \(5y=-5x-15\)，所以直线方程为 \(y=-x-3\).
\item 由斜率定义，\(K=\frac{-5-0}{2-\left(-3\right)}=-1\)，所以直线的斜率为 \(K=-1\).
\item 由直线的斜率定义，直线与 \(OX\) 轴正方向的倾角 \(\alpha\) 满足 \(\tan\alpha=K=-1\)，且 \(0^\circ\leq\alpha<180^\circ\). 由于直线斜率为负，倾角为钝角，故 \(\alpha=135^\circ\).
\end{enumerate}
\end{solution}
